% =====================================================================
\section{Thiết kế thực nghiệm và chỉ số}
\label{sec:experiments}
% =====================================================================

\subsection{Môi trường thực nghiệm}
\label{sec:env}

Toàn bộ thí nghiệm được chạy trong một môi trường ảo Python tách biệt nhằm bảo đảm khả năng
tái lập. Thông tin môi trường dưới đây được đọc trực tiếp từ chính phiên bản phần mềm đã dùng
để huấn luyện và đo thời gian, không suy ra từ phiên bản của trang tài liệu.

\begin{table}[htbp]
  \centering
  \small
  \begin{tabular}{@{}l l@{}}
    \toprule
    Thành phần & Giá trị thực tế \\
    \midrule
    Python & \REnvPython \\
    PyTorch & \REnvTorch \\
    TorchVision & \REnvTorchvision \\
    Bản dựng CUDA & \REnvCudaBuild \\
    GPU & \REnvGpu \\
    VRAM & \REnvVram{} GiB \\
    CPU & \REnvCpu \\
    Số luồng CPU & \REnvCores \\
    RAM & \REnvRam{} GiB \\
    \bottomrule
  \end{tabular}
  \caption{Môi trường thực nghiệm được ghi lại tự động vào \code{results/environment.json} và
  \code{results/self\_test.json}. Số luồng CPU được ghim bằng \code{torch.set\_num\_threads}
  trước các phép đo độ trễ.}
  \label{tab:env}
\end{table}

\paragraph{Tạo môi trường.}
\begin{lstlisting}[language=bash,caption={Cài đặt môi trường ảo và các gói phụ thuộc.},label={lst:setup}]
python -m venv .venv
.venv\Scripts\python.exe -m pip install --upgrade pip
.venv\Scripts\python.exe -m pip install --index-url https://download.pytorch.org/whl/cu128 `
    torch==2.11.0 torchvision==0.26.0
.venv\Scripts\python.exe -m pip install -r requirements.txt
\end{lstlisting}

\subsection{Cấu hình huấn luyện}
\label{sec:train-config}

Để thống nhất số bước cập nhật có giám sát giữa các cấu hình, bảng kết quả chính
\textbf{không dùng Early Stopping}. Cả ba cấu hình chạy đúng 100 epoch và đều được phép chọn
checkpoint bằng cùng một tiêu chí validation. Early Stopping với \code{patience=10} của bản
đề cương ban đầu chỉ có thể dùng cho các thử nghiệm thăm dò và phải được ghi rõ; nó
\textbf{không} được trộn vào bảng so sánh chính vì sẽ tạo ra số epoch khác nhau giữa các cấu
hình.

% Sinh tự động — scripts/make_report_data.py
\begin{table}[htbp]
  \centering
  \small
  \begin{tabular}{@{}l l@{}}
    \toprule
    Thành phần & Thiết lập cho thí nghiệm chính \\
    \midrule
    Optimizer & AdamW \\
    Learning rate & $3 \times 10^{-4}$ \\
    Weight decay & $10^{-4}$ \\
    Betas & $(0{,}9;\ 0{,}999)$ \\
    Eps & $10^{-8}$ \\
    Kích thước minibatch $B_L$ & tối đa 128 ảnh \\
    Kích thước minibatch $B_R$ & bằng $B_L$ tại mỗi bước \\
    Ngân sách học & 100 epoch ($1$ epoch $=$ một lượt duyệt $\DL$) \\
    Scheduler & CosineAnnealingLR, $T_{\max}=100$, $\eta_{\min}=0$ \\
    Hệ số $\lamrot$ & $0{,}5$ cho thí nghiệm chính \\
    Chọn checkpoint & Macro-F1 validation cao nhất; hòa thì loss thấp hơn, rồi epoch sớm hơn \\
    Early stopping & Không dùng trong bảng kết quả chính \\
    Seed & chia split 2026; lần lặp 42, 43, 44 \\
    Chuẩn hóa & mean $=$ std $= (0{,}5;\ 0{,}5;\ 0{,}5)$ \\
    \bottomrule
  \end{tabular}
  \caption{Cấu hình huấn luyện dùng chung cho cả ba cấu hình. Ba cấu hình ghép cặp dùng cùng trọng số khởi tạo của encoder và Superclass Head.}
  \label{tab:hyperparams}
\end{table}

\paragraph{Một epoch được định nghĩa thế nào.}
Một epoch là \emph{một lượt duyệt toàn bộ $\DL$}, không phải một lượt duyệt $\Dtrain$. Nhờ định
nghĩa này, ba mức nhãn có số epoch giống nhau và số bước cập nhật có giám sát tỷ lệ với
$|\DL|$. Minibatch cuối nhỏ hơn 128 được \emph{giữ lại} chứ không bỏ mẫu, nên không có ảnh nào
bị loại khỏi huấn luyện.

\paragraph{Chọn checkpoint.} Tiêu chí là Macro-F1 validation cao nhất. Khi hai epoch có cùng
Macro-F1 (trong sai số $10^{-12}$), ưu tiên loss phân loại validation thấp hơn; nếu vẫn bằng
nhau, chọn epoch sớm hơn. Tập test không tham gia vào bước này dưới bất kỳ hình thức nào.

\subsection{Quy trình và nguyên tắc so sánh}
\label{sec:exp-compare}

Trong mỗi epoch, duyệt $\DL$ một lần với \emph{cùng chuỗi minibatch có giám sát} cho các cấu
hình ghép cặp: cùng seed cho \code{DataLoader}, cùng thứ tự mẫu, cùng phép augmentation. Với mỗi
bước của M8, lấy $B_R$ ngẫu nhiên từ nguồn ảnh quy định tại \tabref{tab:configs}; dùng một bộ
sinh số ngẫu nhiên \textbf{độc lập} để việc lấy ảnh và góc xoay không làm thay đổi augmentation
của nhánh phân loại. Sau đó tính tổng loss và thực hiện một bước optimizer.

\begin{keybox}
\textbf{Vì sao cần bộ sinh số ngẫu nhiên độc lập?} Nếu nhánh phụ dùng chung luồng RNG với nhánh
chính, việc thêm nhánh phụ sẽ \emph{dịch chuyển} chuỗi số ngẫu nhiên của nhánh chính. Khi đó
hai cấu hình ghép cặp sẽ thấy các phép augmentation khác nhau, và chênh lệch $\Delta_s$ trở
thành nhiễu ngẫu nhiên thay vì tín hiệu cần đo.
\end{keybox}

\paragraph{\cfg{E1} không chạy nhánh rotation, kể cả forward.}
Đặt $\lamrot = 0$ nhưng vẫn đưa ảnh xoay qua BatchNorm \textbf{không} được xem là baseline
\cfg{E1}, vì thống kê chạy của mạng vẫn có thể bị thay đổi: các lớp BatchNorm cập nhật
running mean và running variance ngay cả khi gradient của nhánh đó bị nhân với $0$. Trong mã
nguồn, \cfg{E1} được khởi tạo \emph{không có} Rotation Head và vòng huấn luyện bỏ hoàn toàn
nhánh phụ.

\paragraph{Không dùng nhãn của $\DU$.}
Với mọi cấu hình, nhãn của $\DU$ không được dùng để phân tầng minibatch huấn luyện, không được
dùng để xây dựng mục tiêu học, và không được dùng cho bất kỳ quyết định thiết kế nào.

\paragraph{Quy mô khối thực nghiệm.}
Khối chính gồm $3 \text{ cấu hình} \times 3 \text{ mức nhãn} \times 3 \text{ lần lặp} =
\RRuns$ lượt huấn luyện. Ba lần lặp thay đổi \emph{cả} tập con có nhãn \emph{lẫn} ngẫu nhiên
huấn luyện (khởi tạo, thứ tự dữ liệu, augmentation), trong khi phép chia
train/validation/test được cố định. Đây \textbf{không} phải cross-validation: cùng một tập
validation được dùng cho mọi lần lặp, nên ba lần lặp chỉ đo độ ổn định trước hai nguồn ngẫu
nhiên chứ không ước lượng được phương sai do chính phép chia dữ liệu gây ra.

\paragraph{Khảo sát tùy chọn.} $\lamrot \in \{0{,}1;\ 0{,}25;\ 0{,}5;\ 1{,}0\}$ tại mức
$10\,\%$ nhãn, được ưu tiên trước các mở rộng E3--E6. Giữ $\lamrot = 0{,}5$ cho so sánh chính.
Không được lựa chọn $\lam$ hay cấu hình dựa trên kết quả test. Nếu tài nguyên không đủ, phải
điều chỉnh ngân sách chung và ghi nhận \emph{trước khi} mở kết quả test.

\paragraph{Quy trình khóa siêu tham số của phần mở rộng.}
Khối mở rộng contrastive \cfg{E3}--\cfg{E6} có thêm ba siêu tham số phải khóa trước khi đánh giá
test: nhiệt độ SupCon ($0{,}1$), nhiệt độ NT-Xent ($0{,}2$) và hệ số $\lamc$. Hai nhiệt độ và
chiều Projection Head ($128$) được chốt theo giá trị chuẩn của tài liệu
gốc~\cite{chen2020simclr,khosla2020supcon}; riêng $\lamc$ được chọn bằng \textbf{Macro-F1
validation} qua một script \emph{không có} bước đánh giá test (\secref{sec:results-ext-lambda}).
Quy tắc chọn được khai báo trước khi chạy để tránh điều chỉnh giao thức sau khi đã thấy kết quả.

\subsection{Chỉ số và cách báo cáo}
\label{sec:metrics}

Chỉ số chính là \textbf{Macro-F1} của Animal và Vehicle. Báo cáo thêm Accuracy,
precision/recall/F1 từng lớp và ma trận nhầm lẫn trên toàn bộ 10.000 ảnh test. Mỗi ảnh chỉ nhận
một dự đoán của nhánh siêu lớp; \textbf{không} dùng test-time augmentation trong thí nghiệm
chính.

Các chỉ số được định nghĩa như sau:

\begin{align}
  \mathrm{Accuracy} &= \frac{1}{N}\sum_{i=1}^{N} \mathbb{1}\!\left[\hat{y}_i = y_i\right],
  \label{eq:acc}\\
  P_c &= \frac{\mathrm{TP}_c}{\mathrm{TP}_c + \mathrm{FP}_c}, \qquad
  R_c = \frac{\mathrm{TP}_c}{\mathrm{TP}_c + \mathrm{FN}_c}, \label{eq:pr}\\
  \mathrm{F1}_c &= \frac{2 P_c R_c}{P_c + R_c}, \qquad
  \mathrm{Macro\text{-}F1} = \frac{\mathrm{F1}_{\mathrm{Animal}} + \mathrm{F1}_{\mathrm{Vehicle}}}{2}.
  \label{eq:macro}
\end{align}

\paragraph{Quy ước biên.} Nếu mẫu số của precision, recall hoặc F1 bằng $0$, giá trị tương ứng
được quy ước bằng $0$ và dùng thống nhất cho mọi cấu hình. Chỉ số được tính từ \emph{toàn bộ}
dự đoán của tập đánh giá; không lấy trung bình không trọng số các F1 của từng minibatch, vì
làm như vậy sẽ cho trọng số sai cho các minibatch nhỏ ở cuối epoch.

\paragraph{Cách tổng hợp.} Báo cáo trung bình $\pm$ độ lệch chuẩn mẫu qua ba lần lặp và chênh
lệch ghép cặp $\Delta_s = \mathrm{Macro\text{-}F1}_s(\text{M8}) -
\mathrm{Macro\text{-}F1}_s(\cfg{E1})$ tính bằng điểm phần trăm. Mức $10\,\%$ là so sánh ưu
tiên; mức $20\,\%$ và $50\,\%$ dùng để đánh giá xu hướng.

\begin{warnbox}
\textbf{Giới hạn suy luận thống kê.} Ba lần lặp chỉ cho bằng chứng ban đầu về độ ổn định.
Với $n=3$ và độ lệch chuẩn quan sát được, khoảng tin cậy rất rộng và không có cơ sở để tuyên
bố \enquote{có ý nghĩa thống kê} hay \enquote{luôn tốt hơn}. Đặc biệt, \textbf{không} được chọn
một lần chạy tốt nhất rồi kết luận từ nó: báo cáo dùng trung bình và độ lệch chuẩn qua cả ba
lần~\cite{bouthillier2021variance}.
\end{warnbox}

\paragraph{Rotation Accuracy.} Chỉ là chỉ số \emph{chẩn đoán} cho nhánh phụ. Nếu đo trên
validation, phải đánh giá đủ bốn phép xoay cho mỗi ảnh và ghi rõ giao thức. Độ chính xác cao ở
nhiệm vụ phụ \textbf{không} thay thế bằng chứng cải thiện Macro-F1 của nhiệm vụ chính; một
nhiệm vụ phụ được học tốt vẫn hoàn toàn có thể không giúp gì cho nhiệm vụ chính.

\subsection{Chi phí suy luận và giới hạn kết luận về IoT}
\label{sec:cost-protocol}

\paragraph{Đếm tham số và dung lượng.} Đếm riêng số tham số khi huấn luyện và khi suy luận; đo
dung lượng tệp chỉ chứa trọng số suy luận, \textbf{không} gộp optimizer state (vì optimizer
state không tồn tại trên thiết bị triển khai). Với cấu hình tại \secref{sec:arch}, mô hình suy
luận có \RParamInference{} tham số; Rotation Head thêm \RParamRotationHead{} tham số. Phần tham
số FP32 của mô hình suy luận tương đương khoảng \RParamInferenceMiB{} MiB, \emph{chưa} gồm
buffer, activation và bộ nhớ của runtime~\cite{pytorch_resnet}. Đây là thông số kiến trúc, không
phải kết quả huấn luyện.

\paragraph{Giao thức đo độ trễ.} Đo với batch $=1$, tensor $1 \times 3 \times 32 \times 32$,
FP32, \code{model.eval()} và \code{torch.inference\_mode()}. Dùng \emph{cùng} máy, cùng runtime
và cùng số luồng CPU cho mọi cấu hình. Chạy khởi động 50 lượt, sau đó đo tối thiểu 1.000 lượt
trên cùng các ảnh validation đã tiền xử lý. Nếu đo trên GPU, đồng bộ trước và sau khối tính thời
gian. Báo cáo mean, median, p95 theo ms/ảnh; lặp ba phiên và lưu log~\cite{pytorch_benchmark}.

\paragraph{Tách riêng các thành phần thời gian.} Phép đo trên chỉ tính thời gian \emph{suy luận
của mô hình}. Thời gian đọc ảnh, tiền xử lý và truyền dữ liệu phải được đo và báo cáo riêng.
Hai mô hình \cfg{E1} và \cfg{E2} có \emph{cùng} kiến trúc suy luận, nên không kỳ vọng giảm độ
phức tạp kiến trúc chỉ nhờ bỏ head phụ.
